%%%%%%%%%%%%%%%%%%%%%%%%%%%%%%%%%%%%%%%%%%%%%%%%%%%%%%%%%%%%%
%% THÉORIE DES CATÉGORIES — Chapitres 5–6
%% Lemme de Yoneda et Représentabilité / Monades et Algèbres
%%%%%%%%%%%%%%%%%%%%%%%%%%%%%%%%%%%%%%%%%%%%%%%%%%%%%%%%%%%%%

%%%%%%%%%%%%%%%%%%%%%%%%%%%%%%%%%%%%%%%%%%%%%%%%%%%%%%%%%%%%
%%     CHAPITRE 5 : LEMME DE YONEDA ET REPRÉSENTABILITÉ   %%
%%%%%%%%%%%%%%%%%%%%%%%%%%%%%%%%%%%%%%%%%%%%%%%%%%%%%%%%%%%%

\chapter{Lemme de Yoneda et représentabilité}\label{ch:yoneda}
\minitoc

\vspace{1em}

Le lemme de Yoneda\index{Yoneda, lemme de|(} est, selon beaucoup de catégoriciens,
le résultat le plus fondamental de la théorie des catégories.
Il affirme qu'un objet est entièrement déterminé
par l'ensemble des morphismes qui y arrivent (ou qui en partent),
et que cette détermination est fonctorielle.
Plus précisément, le lemme établit une bijection naturelle entre les
transformations naturelles issues d'un foncteur représentable et les
éléments de l'objet cible, ce qui permet de \emph{plonger} toute catégorie
dans une catégorie de foncteurs.

Ce chapitre développe la théorie des préfaisceaux, démontre le lemme de Yoneda
dans sa forme complète, en déduit le plongement de Yoneda, puis étudie
les foncteurs représentables et leurs propriétés.

\index{préfaisceau|(}
\index{foncteur représentable|(}

%% ============================================================
\section{Préfaisceaux et foncteurs \texorpdfstring{$\Hom$}{Hom}}\label{sec:prefaisceaux}
%% ============================================================

\begin{definition}[Préfaisceau]\label{def:prefaisceau}
\index{préfaisceau!définition}
\index{catégorie!des préfaisceaux}
Soit $\mathcal{C}$ une catégorie localement petite.
Un \textbf{préfaisceau} sur $\mathcal{C}$ est un foncteur
\[
    F \colon \mathcal{C}^{\op} \longrightarrow \Set.
\]
La catégorie des préfaisceaux sur $\mathcal{C}$ est la catégorie de foncteurs
$\Fun(\mathcal{C}^{\op}, \Set)$,
que l'on note $\widehat{\mathcal{C}}$ ou $\mathrm{PSh}(\mathcal{C})$.
\index{$\widehat{\mathcal{C}}$}
\index{PSh@$\mathrm{PSh}(\mathcal{C})$}
Les morphismes dans $\widehat{\mathcal{C}}$ sont les transformations naturelles.
\end{definition}

\begin{remarque}\label{rem:prefaisceau-taille}
\index{préfaisceau!questions de taille}
La catégorie $\widehat{\mathcal{C}}$ n'est en général pas petite, même si
$\mathcal{C}$ est petite : les foncteurs de $\mathcal{C}^{\op}$ vers $\Set$
forment une classe propre dès que $\mathcal{C}$ a suffisamment d'objets.
Toutefois, si $\mathcal{C}$ est petite, alors $\widehat{\mathcal{C}}$ est
localement petite et admet toutes les limites et colimites petites
(calculées argument par argument).
\end{remarque}

\begin{definition}[Foncteur $\Hom$ covariant]\label{def:hom-covariant}
\index{foncteur Hom@foncteur $\Hom$!covariant}
\index{h@$h^A$}
Soit $\mathcal{C}$ une catégorie localement petite et $A \in \Ob(\mathcal{C})$.
Le \textbf{foncteur $\Hom$ covariant} représenté par $A$ est le foncteur
\[
    h^A = \Hom(A, -) \colon \mathcal{C} \longrightarrow \Set
\]
défini sur les objets par $h^A(B) = \Hom(A, B)$ et sur les morphismes
$f \colon B \to C$ par la post-composition :
\[
    h^A(f) = f_* \colon \Hom(A, B) \longrightarrow \Hom(A, C), \quad
    g \longmapsto f \circ g.
\]
\end{definition}

\begin{definition}[Foncteur $\Hom$ contravariant]\label{def:hom-contravariant}
\index{foncteur Hom@foncteur $\Hom$!contravariant}
\index{h@$h_A$}
Soit $A \in \Ob(\mathcal{C})$.
Le \textbf{foncteur $\Hom$ contravariant} représenté par $A$ est le foncteur
\[
    h_A = \Hom(-, A) \colon \mathcal{C}^{\op} \longrightarrow \Set
\]
défini sur les objets par $h_A(B) = \Hom(B, A)$ et sur les morphismes
$f \colon B \to C$ (dans $\mathcal{C}$, donc $f \colon C \to B$ dans
$\mathcal{C}^{\op}$) par la pré-composition :
\[
    h_A(f) = f^* \colon \Hom(C, A) \longrightarrow \Hom(B, A), \quad
    g \longmapsto g \circ f.
\]
Ainsi $h_A$ est un préfaisceau sur $\mathcal{C}$.
\end{definition}

\begin{exemple}\label{ex:hom-foncteurs}
\index{foncteur Hom@foncteur $\Hom$!exemples}
\begin{enumerate}[label=(\roman*)]
\item Dans $\Set$, le foncteur $h^{\{*\}} = \Hom(\{*\}, -)$ est
    (naturellement isomorphe à) le foncteur identité $\Id_{\Set}$,
    car un élément de $\Hom(\{*\}, X)$ correspond exactement
    à un élément de $X$.

\item Dans $\Ab$, le foncteur $h^{\Z} = \Hom(\Z, -)$ est le
    foncteur oubli\index{foncteur oubli} $U \colon \Ab \to \Set$,
    car un morphisme de groupes $\Z \to G$ est déterminé par l'image de $1$.

\item Dans $\Top$, le foncteur $h^{\{*\}}$ est le foncteur
    ensemble sous-jacent. Le foncteur $h^{[0,1]}$ est le foncteur
    des chemins\index{foncteur!des chemins}, dont l'étude mène à la
    théorie de l'homotopie.

\item Dans $\mathsf{Vect}_k$, le foncteur $h_V = \Hom(-, V)$
    envoie $W$ sur l'espace des applications linéaires $W \to V$.
    En particulier, $h_{k} = \Hom(-, k)$ est le foncteur dual\index{foncteur dual}.
\end{enumerate}
\end{exemple}

\begin{definition}[Foncteur $\Hom$ bifunctoriel]\label{def:hom-bifuncteur}
\index{foncteur Hom@foncteur $\Hom$!bifuncteur}
\index{bifuncteur}
Le \textbf{foncteur $\Hom$ bifunctoriel} est le foncteur
\[
    \Hom(-,-) \colon \mathcal{C}^{\op} \times \mathcal{C} \longrightarrow \Set
\]
défini sur les paires d'objets par $\Hom(A, B)$ et sur les paires de morphismes
$(f \colon A' \to A,\, g \colon B \to B')$ par
\[
    \Hom(f, g) \colon \Hom(A, B) \longrightarrow \Hom(A', B'), \quad
    h \longmapsto g \circ h \circ f.
\]
\end{definition}

\begin{proposition}\label{prop:hom-bifuncteur-fonctorialite}
\index{foncteur Hom@foncteur $\Hom$!fonctorialité}
Le foncteur $\Hom(-,-)$ est bien un bifuncteur : pour tous morphismes
composables $(f_1, g_1)$ et $(f_2, g_2)$, on a
\[
    \Hom(f_1 \circ f_2,\, g_2 \circ g_1) = \Hom(f_2, g_2) \circ \Hom(f_1, g_1),
\]
et $\Hom(\id_A, \id_B) = \id_{\Hom(A,B)}$.
\end{proposition}

\begin{proof}
Pour tout $h \in \Hom(A, B)$, on a
\begin{align*}
    \Hom(f_1 \circ f_2,\, g_2 \circ g_1)(h)
    &= (g_2 \circ g_1) \circ h \circ (f_1 \circ f_2) \\
    &= g_2 \circ (g_1 \circ h \circ f_1) \circ f_2 \\
    &= \Hom(f_2, g_2)\bigl(\Hom(f_1, g_1)(h)\bigr).
\end{align*}
L'assertion sur les identités est immédiate.
\end{proof}


%% ============================================================
\section{Le plongement de Yoneda}\label{sec:plongement-yoneda}
%% ============================================================

\begin{definition}[Plongement de Yoneda covariant]\label{def:plongement-yoneda-cov}
\index{plongement de Yoneda!covariant}
\index{Yoneda!plongement}
Le \textbf{plongement de Yoneda covariant} est le foncteur
\[
    \mathbf{y} \colon \mathcal{C} \longrightarrow \Fun(\mathcal{C}^{\op}, \Set)
    = \widehat{\mathcal{C}}, \qquad
    A \longmapsto h_A = \Hom(-, A).
\]
Sur les morphismes $f \colon A \to B$, il est défini par la transformation
naturelle $\mathbf{y}(f) = f_* \colon h_A \Rightarrow h_B$ dont la composante
en $C \in \mathcal{C}$ est
\[
    (f_*)_C \colon \Hom(C, A) \longrightarrow \Hom(C, B), \qquad
    g \longmapsto f \circ g.
\]
\end{definition}

\begin{definition}[Plongement de Yoneda contravariant]\label{def:plongement-yoneda-contrav}
\index{plongement de Yoneda!contravariant}
De même, le \textbf{plongement de Yoneda contravariant} est le foncteur
\[
    \mathbf{y}^{\op} \colon \mathcal{C}^{\op} \longrightarrow
    \Fun(\mathcal{C}, \Set), \qquad
    A \longmapsto h^A = \Hom(A, -).
\]
\end{definition}

\begin{remarque}\label{rem:yoneda-notations}
\index{plongement de Yoneda!notations}
Suivant les auteurs, le plongement de Yoneda covariant est parfois noté
$\mathcal{C} \hookrightarrow \widehat{\mathcal{C}}$,
$A \mapsto \mathcal{C}(-, A)$, ou encore $A \mapsto \mathrm{Yo}(A)$.
Nous utiliserons principalement la notation $h_A$ pour le foncteur
contravariant représenté par $A$ et $h^A$ pour le covariant.
\end{remarque}

\begin{lemme}\label{lem:y-bien-defini}
\index{plongement de Yoneda!bonne définition}
Le plongement de Yoneda $\mathbf{y}$ est un foncteur bien défini :
pour tout morphisme $f \colon A \to B$ dans $\mathcal{C}$,
la famille $(f_*)_C = f \circ (-)$ est une transformation naturelle
$h_A \Rightarrow h_B$.
\end{lemme}

\begin{proof}
Il faut vérifier que pour tout morphisme $u \colon C' \to C$ dans $\mathcal{C}$,
le diagramme suivant commute :
\[
\begin{tikzcd}
\Hom(C, A) \arrow[r, "(f_*)_C"] \arrow[d, "u^*"'] &
\Hom(C, B) \arrow[d, "u^*"] \\
\Hom(C', A) \arrow[r, "(f_*)_{C'}"'] &
\Hom(C', B)
\end{tikzcd}
\]
Pour tout $g \in \Hom(C, A)$, on a
\[
    u^* \circ (f_*)_C(g) = u^*(f \circ g) = f \circ g \circ u
    = (f_*)_{C'}(g \circ u) = (f_*)_{C'} \circ u^*(g).
\]
Donc le diagramme commute.
\end{proof}


%% ============================================================
\section{Le lemme de Yoneda}\label{sec:lemme-yoneda}
%% ============================================================

Nous énonçons et démontrons maintenant le résultat central de ce chapitre.
Nous traitons d'abord la version contravariante (préfaisceaux), puis
indiquons la version covariante.

\begin{theoreme}[Lemme de Yoneda]\label{thm:lemme-yoneda}
\index{Yoneda, lemme de!énoncé}
\index{transformation naturelle!et lemme de Yoneda}
Soit $\mathcal{C}$ une catégorie localement petite,
$A \in \Ob(\mathcal{C})$, et $F \colon \mathcal{C}^{\op} \to \Set$
un préfaisceau. Il existe une bijection
\[
    \Phi_A \colon \Nat(h_A, F) \xrightarrow{\;\sim\;} F(A)
\]
qui est naturelle en $A$ et en $F$.
Explicitement, $\Phi_A$ est définie par
\[
    \Phi_A(\alpha) = \alpha_A(\id_A)
\]
pour toute transformation naturelle $\alpha \colon h_A \Rightarrow F$.
\end{theoreme}

\begin{proof}
\index{Yoneda, lemme de!preuve}
La preuve se décompose en quatre étapes : construction de la bijection,
construction de l'inverse, vérification que ces constructions sont
mutuellement inverses, et preuve de la naturalité.

\medskip
\textbf{Étape 1 : Définition de $\Phi_A$.}
Soit $\alpha \colon h_A \Rightarrow F$ une transformation naturelle.
Pour chaque objet $B$ de $\mathcal{C}$, $\alpha$ possède une composante
$\alpha_B \colon \Hom(B, A) \to F(B)$.
En particulier, pour $B = A$, on obtient
$\alpha_A \colon \Hom(A, A) \to F(A)$.
On pose
\[
    \Phi_A(\alpha) \coloneqq \alpha_A(\id_A) \in F(A).
\]

\medskip
\textbf{Étape 2 : Construction de l'inverse $\Psi_A$.}
Soit $x \in F(A)$. On définit une transformation naturelle
$\Psi_A(x) = \beta^x \colon h_A \Rightarrow F$ en posant,
pour chaque objet $B$ de $\mathcal{C}$ et chaque morphisme
$f \in \Hom(B, A) = h_A(B)$,
\[
    \beta^x_B(f) \coloneqq F(f)(x) \in F(B).
\]
Ici, $F(f) \colon F(A) \to F(B)$ est l'application induite par $F$
(rappelons que $F$ est un foncteur de $\mathcal{C}^{\op}$ vers $\Set$,
donc un morphisme $f \colon B \to A$ dans $\mathcal{C}$ induit
$F(f) \colon F(A) \to F(B)$).

Vérifions que $\beta^x$ est bien une transformation naturelle.
Pour un morphisme $u \colon C \to B$ dans $\mathcal{C}$,
il faut montrer que le diagramme suivant commute :
\[
\begin{tikzcd}
\Hom(B, A) \arrow[r, "\beta^x_B"] \arrow[d, "u^*"'] &
F(B) \arrow[d, "F(u)"] \\
\Hom(C, A) \arrow[r, "\beta^x_C"'] &
F(C)
\end{tikzcd}
\]
Pour tout $f \in \Hom(B, A)$, on a
\begin{align*}
    F(u) \circ \beta^x_B(f)
    &= F(u)\bigl(F(f)(x)\bigr) \\
    &= \bigl(F(u) \circ F(f)\bigr)(x) \\
    &= F(f \circ u)(x) \quad \text{(fonctorialité de $F$ sur $\mathcal{C}^{\op}$)} \\
    &= \beta^x_C(f \circ u) \\
    &= \beta^x_C(u^*(f)).
\end{align*}
Donc $\beta^x$ est bien naturelle.

\medskip
\textbf{Étape 3 : $\Phi_A$ et $\Psi_A$ sont mutuellement inverses.}

\emph{Vérification $\Phi_A \circ \Psi_A = \id_{F(A)}$.}
Pour tout $x \in F(A)$,
\[
    \Phi_A(\Psi_A(x))
    = \Phi_A(\beta^x)
    = \beta^x_A(\id_A)
    = F(\id_A)(x)
    = \id_{F(A)}(x)
    = x.
\]

\emph{Vérification $\Psi_A \circ \Phi_A = \id_{\Nat(h_A, F)}$.}
Pour toute transformation naturelle $\alpha \colon h_A \Rightarrow F$,
posons $x = \Phi_A(\alpha) = \alpha_A(\id_A)$.
Il faut montrer que $\beta^x = \alpha$,
c'est-à-dire que pour tout objet $B$ et tout $f \in \Hom(B, A)$,
\[
    \beta^x_B(f) = \alpha_B(f).
\]
Par définition, $\beta^x_B(f) = F(f)(x) = F(f)(\alpha_A(\id_A))$.
Or la naturalité de $\alpha$ appliquée au morphisme $f \colon B \to A$
donne le diagramme commutatif :
\[
\begin{tikzcd}
\Hom(A, A) \arrow[r, "\alpha_A"] \arrow[d, "f^*"'] &
F(A) \arrow[d, "F(f)"] \\
\Hom(B, A) \arrow[r, "\alpha_B"'] &
F(B)
\end{tikzcd}
\]
En évaluant sur $\id_A \in \Hom(A, A)$, on obtient
\[
    F(f)(\alpha_A(\id_A)) = \alpha_B(f^*(\id_A)) = \alpha_B(\id_A \circ f)
    = \alpha_B(f).
\]
Donc $\beta^x_B(f) = \alpha_B(f)$ pour tout $B$ et tout $f$,
ce qui montre $\Psi_A(\Phi_A(\alpha)) = \alpha$.

\medskip
\textbf{Étape 4 : Naturalité.}

\emph{Naturalité en $A$.}
Soit $u \colon A' \to A$ un morphisme dans $\mathcal{C}$.
Il faut vérifier que le diagramme suivant commute :
\[
\begin{tikzcd}
\Nat(h_A, F) \arrow[r, "\Phi_A"] \arrow[d, "u^*"'] &
F(A) \arrow[d, "F(u)"] \\
\Nat(h_{A'}, F) \arrow[r, "\Phi_{A'}"'] &
F(A')
\end{tikzcd}
\]
où $u^*(\alpha) = \alpha \circ \mathbf{y}(u) = \alpha \circ u_*$,
c'est-à-dire la transformation naturelle dont la composante en $B$ est
$\alpha_B \circ (u \circ -)$.
Pour toute $\alpha \colon h_A \Rightarrow F$,
\begin{align*}
    F(u)(\Phi_A(\alpha))
    &= F(u)(\alpha_A(\id_A)), \\
    \Phi_{A'}(u^*(\alpha))
    &= (u^*(\alpha))_{A'}(\id_{A'})
    = \alpha_{A'}(u \circ \id_{A'})
    = \alpha_{A'}(u).
\end{align*}
Par naturalité de $\alpha$ en le morphisme $u \colon A' \to A$,
$F(u)(\alpha_A(\id_A)) = \alpha_{A'}(u)$. Le diagramme commute.

\emph{Naturalité en $F$.}
Soit $\eta \colon F \Rightarrow G$ une transformation naturelle
entre préfaisceaux. On vérifie que
\[
\begin{tikzcd}
\Nat(h_A, F) \arrow[r, "\Phi^F_A"] \arrow[d, "\eta_*"'] &
F(A) \arrow[d, "\eta_A"] \\
\Nat(h_A, G) \arrow[r, "\Phi^G_A"'] &
G(A)
\end{tikzcd}
\]
commute, où $\eta_*(\alpha) = \eta \circ \alpha$. Pour toute $\alpha$,
\[
    \eta_A(\Phi^F_A(\alpha)) = \eta_A(\alpha_A(\id_A))
    = (\eta \circ \alpha)_A(\id_A)
    = \Phi^G_A(\eta_*(\alpha)). \qedhere
\]
\end{proof}

\begin{remarque}\label{rem:yoneda-covariant}
\index{Yoneda, lemme de!version covariante}
La version \emph{covariante} du lemme de Yoneda s'énonce comme suit.
Pour tout foncteur $G \colon \mathcal{C} \to \Set$ et tout
$A \in \Ob(\mathcal{C})$, on a une bijection naturelle
\[
    \Nat(h^A, G) \cong G(A),
    \qquad \alpha \mapsto \alpha_A(\id_A).
\]
La preuve est entièrement analogue.
\end{remarque}


%% ============================================================
\section{Corollaires du lemme de Yoneda}\label{sec:yoneda-corollaires}
%% ============================================================

\begin{corollaire}[Le plongement de Yoneda est pleinement fidèle]%
\label{cor:yoneda-pleinement-fidele}
\index{plongement de Yoneda!pleinement fidèle}
\index{foncteur!pleinement fidèle}
Le foncteur $\mathbf{y} \colon \mathcal{C} \to \widehat{\mathcal{C}}$,
$A \mapsto h_A$, est pleinement fidèle.
Autrement dit, pour tous objets $A, B$ de $\mathcal{C}$,
\[
    \mathbf{y}_{A,B} \colon \Hom_{\mathcal{C}}(A, B)
    \xrightarrow{\;\sim\;} \Nat(h_A, h_B)
\]
est une bijection.
\end{corollaire}

\begin{proof}
Appliquons le lemme de Yoneda~\ref{thm:lemme-yoneda} au préfaisceau
$F = h_B = \Hom(-, B)$. On obtient la bijection naturelle
\[
    \Phi_A \colon \Nat(h_A, h_B) \xrightarrow{\;\sim\;} h_B(A) = \Hom(A, B).
\]
Par construction, $\Phi_A(\alpha) = \alpha_A(\id_A)$ pour toute
transformation naturelle $\alpha \colon h_A \Rightarrow h_B$.

Or, par définition du plongement de Yoneda, pour tout morphisme
$f \colon A \to B$, on a $\mathbf{y}(f) = f_*$ avec
$(f_*)_C(g) = f \circ g$ pour tout $g \in \Hom(C, A)$.
Donc
\[
    \Phi_A(\mathbf{y}(f)) = (f_*)_A(\id_A) = f \circ \id_A = f.
\]
Cela montre que $\Phi_A \circ \mathbf{y}_{A,B} = \id_{\Hom(A,B)}$,
donc $\mathbf{y}_{A,B}$ est l'inverse de $\Phi_A$, et en particulier
c'est une bijection.
\end{proof}

\begin{corollaire}\label{cor:h-iso-implique-iso}
\index{Yoneda, lemme de!conséquence isomorphisme}
\index{isomorphisme!et foncteurs représentables}
Soient $A, B$ des objets d'une catégorie localement petite $\mathcal{C}$.
Alors
\[
    h_A \cong h_B \quad \text{dans } \widehat{\mathcal{C}}
    \qquad \Longleftrightarrow \qquad
    A \cong B \quad \text{dans } \mathcal{C}.
\]
De même, $h^A \cong h^B$ dans $\Fun(\mathcal{C}, \Set)$
si et seulement si $A \cong B$ dans $\mathcal{C}$.
\end{corollaire}

\begin{proof}
$(\Leftarrow)$ Si $f \colon A \xrightarrow{\sim} B$ est un isomorphisme,
alors $\mathbf{y}(f) = f_* \colon h_A \Rightarrow h_B$ est un isomorphisme
dans $\widehat{\mathcal{C}}$ car $\mathbf{y}$ est un foncteur.

$(\Rightarrow)$ Supposons $\alpha \colon h_A \xrightarrow{\sim} h_B$
un isomorphisme de préfaisceaux.
Puisque $\mathbf{y}$ est pleinement fidèle
(corollaire~\ref{cor:yoneda-pleinement-fidele}), il existe un unique
$f \colon A \to B$ tel que $\mathbf{y}(f) = \alpha$,
et un unique $g \colon B \to A$ tel que $\mathbf{y}(g) = \alpha^{-1}$.
Alors
\[
    \mathbf{y}(g \circ f) = \mathbf{y}(g) \circ \mathbf{y}(f)
    = \alpha^{-1} \circ \alpha = \id_{h_A} = \mathbf{y}(\id_A),
\]
et la fidélité de $\mathbf{y}$ donne $g \circ f = \id_A$.
De même, $f \circ g = \id_B$, donc $f$ est un isomorphisme.
\end{proof}

\begin{corollaire}\label{cor:yoneda-elements}
\index{Yoneda, lemme de!et éléments}
Pour tout préfaisceau $F \in \widehat{\mathcal{C}}$ et tout objet
$A \in \mathcal{C}$, les éléments de $F(A)$ correspondent
bijectivement aux morphismes $h_A \to F$ dans $\widehat{\mathcal{C}}$.
En particulier, les «~éléments généralisés de type $A$~» de $F$
sont exactement les transformations naturelles $h_A \Rightarrow F$.
\end{corollaire}

\begin{remarque}\label{rem:yoneda-philosophie}
\index{Yoneda, lemme de!interprétation philosophique}
Le lemme de Yoneda exprime le principe suivant : \emph{un objet est
entièrement déterminé par ses relations avec tous les autres objets}.
Plus précisément, connaître le préfaisceau $h_A = \Hom(-, A)$ revient
à connaître $A$ à isomorphisme unique près. Ce point de vue
est central en géométrie algébrique, où les «~points à valeurs dans $S$~»
d'un schéma $X$ sont les morphismes $S \to X$, c'est-à-dire les éléments
de $h_X(S)$.
\end{remarque}


%% ============================================================
\section{Foncteurs représentables}\label{sec:foncteurs-representables}
%% ============================================================

\begin{definition}[Foncteur représentable]\label{def:foncteur-representable}
\index{foncteur représentable!définition}
\index{objet représentant}
Un préfaisceau $F \colon \mathcal{C}^{\op} \to \Set$ est dit
\textbf{représentable} s'il existe un objet $A \in \mathcal{C}$
et un isomorphisme naturel $\alpha \colon h_A \xrightarrow{\sim} F$.
L'objet $A$ est appelé l'\textbf{objet représentant} de $F$,
et le couple $(A, \alpha)$ est une \textbf{représentation} de $F$.

Par le lemme de Yoneda, donner un tel isomorphisme revient
à donner un élément $\xi = \alpha_A(\id_A) \in F(A)$
tel que pour tout objet $B$ et tout $x \in F(B)$,
il existe un unique $f \colon B \to A$ avec $F(f)(\xi) = x$.
L'élément $\xi$ est l'\textbf{élément universel}\index{élément universel}.
\end{definition}

\begin{exemple}\label{ex:representables}
\index{foncteur représentable!exemples}
\begin{enumerate}[label=(\roman*)]
\item \textbf{Foncteur oubli des groupes.}\index{foncteur oubli!des groupes}
    Le foncteur oubli $U \colon \Grp \to \Set$ est représenté par $\Z$,
    car $\Hom_{\Grp}(\Z, G) \cong U(G)$ via $\varphi \mapsto \varphi(1)$.
    L'élément universel est $1 \in U(\Z) = \Z$.

\item \textbf{Foncteur ensemble de puissance.}\index{foncteur!ensemble de puissance}
    Le foncteur $\mathcal{P} \colon \Set^{\op} \to \Set$ qui envoie
    $X$ sur $\mathcal{P}(X) = \{A \subseteq X\}$ est représenté
    par l'ensemble $\{0, 1\}$ : un sous-ensemble $A \subseteq X$
    correspond à sa fonction caractéristique $\chi_A \colon X \to \{0,1\}$.

\item \textbf{Produits.}\index{produit!comme foncteur représentable}
    Dans une catégorie $\mathcal{C}$ avec produits binaires, le foncteur
    $\mathcal{C}^{\op} \to \Set$ défini par
    $X \mapsto \Hom(X, A) \times \Hom(X, B)$
    est représenté par $A \times B$.

\item \textbf{Grassmanniennes.}\index{Grassmannienne}
    En géométrie algébrique, le foncteur de Grassmann
    $\mathrm{Grass}(r, n) \colon \mathsf{Sch}^{\op} \to \Set$
    qui envoie un schéma $S$ sur l'ensemble des sous-fibrés localement
    libres de rang $r$ de $\mathcal{O}_S^n$ est représentable par
    la Grassmannienne $\mathrm{Gr}(r, n)$.
\end{enumerate}
\end{exemple}

\begin{proposition}\label{prop:representant-unique}
\index{foncteur représentable!unicité du représentant}
Si $F$ est un préfaisceau représentable, son objet représentant
est unique à isomorphisme unique près.
\end{proposition}

\begin{proof}
Si $(A, \alpha)$ et $(B, \beta)$ sont deux représentations de $F$,
alors $\beta^{-1} \circ \alpha \colon h_A \xrightarrow{\sim} h_B$
est un isomorphisme de préfaisceaux. Par le
corollaire~\ref{cor:h-iso-implique-iso}, $A \cong B$ dans $\mathcal{C}$.
L'unicité de l'isomorphisme découle de la fidélité de $\mathbf{y}$.
\end{proof}


%% ============================================================
\section{Tout préfaisceau est colimite de représentables}%
\label{sec:prefaisceau-colimite}
%% ============================================================

\begin{definition}[Catégorie des éléments]\label{def:categorie-elements}
\index{catégorie des éléments}
\index{$\int_{\mathcal{C}} F$}
Soit $F \colon \mathcal{C}^{\op} \to \Set$ un préfaisceau.
La \textbf{catégorie des éléments} de $F$, notée
$\int_{\mathcal{C}} F$ ou $\mathrm{el}(F)$, est la catégorie dont :
\begin{itemize}
    \item les objets sont les paires $(C, x)$ avec $C \in \mathcal{C}$
    et $x \in F(C)$ ;
    \item un morphisme $(C, x) \to (C', x')$ est un morphisme
    $f \colon C \to C'$ dans $\mathcal{C}$ tel que $F(f)(x') = x$.
\end{itemize}
\end{definition}

\begin{proposition}\label{prop:el-projection}
\index{catégorie des éléments!foncteur projection}
Le foncteur de projection $\pi \colon \int_{\mathcal{C}} F \to \mathcal{C}$
défini par $\pi(C, x) = C$ et $\pi(f) = f$ est un foncteur fidèle.
De plus, la catégorie $\int_{\mathcal{C}} F$ est petite lorsque
$\mathcal{C}$ est petite.
\end{proposition}

\begin{proof}
La fidélité est immédiate puisque $\pi$ est injectif sur les morphismes.
Si $\mathcal{C}$ est petite, alors $\Ob(\int_{\mathcal{C}} F)
= \coprod_{C \in \mathcal{C}} F(C)$ est un ensemble, et de même pour
les morphismes.
\end{proof}

\begin{theoreme}[Tout préfaisceau est colimite de représentables]%
\label{thm:prefaisceau-colimite-representables}
\index{préfaisceau!comme colimite de représentables}
\index{colimite!de représentables}
Soit $\mathcal{C}$ une catégorie petite et
$F \colon \mathcal{C}^{\op} \to \Set$ un préfaisceau. Alors
\[
    F \cong \colim_{(C, x) \in \int_{\mathcal{C}} F} h_C
\]
dans $\widehat{\mathcal{C}}$, où la colimite est prise le long du
foncteur $\int_{\mathcal{C}} F \to \widehat{\mathcal{C}}$,
$(C, x) \mapsto h_C$.
\end{theoreme}

\begin{proof}
Notons $G = \colim_{(C,x) \in \int_{\mathcal{C}} F} h_C$.
Par le lemme de Yoneda, pour tout préfaisceau $P$,
\begin{align*}
    \Nat(G, P)
    &\cong \lim_{(C,x) \in (\int_{\mathcal{C}} F)^{\op}} \Nat(h_C, P) \\
    &\cong \lim_{(C,x) \in (\int_{\mathcal{C}} F)^{\op}} P(C).
\end{align*}
Un élément de cette limite est une famille compatible
$(s_{(C,x)})_{(C,x)} \in \prod_{(C,x)} P(C)$ telle que
pour tout morphisme $f \colon (C,x) \to (C', x')$ dans $\int F$,
on a $P(f)(s_{(C',x')}) = s_{(C,x)}$.
Or une transformation naturelle $\eta \colon F \Rightarrow P$ est
une famille $(\eta_C)_{C \in \mathcal{C}}$ avec
$\eta_C \colon F(C) \to P(C)$ naturelle en $C$.
La correspondance est $\eta \mapsto (\eta_C(x))_{(C,x)}$,
et on vérifie que c'est une bijection. Ainsi
$\Nat(G, P) \cong \Nat(F, P)$ naturellement en $P$,
et le lemme de Yoneda (appliqué dans $\widehat{\mathcal{C}}$)
donne $G \cong F$.
\end{proof}

\begin{remarque}\label{rem:colimite-representables-densite}
\index{Yoneda!densité}
\index{théorème de densité}
Le théorème~\ref{thm:prefaisceau-colimite-representables} est souvent
appelé \emph{théorème de densité} de Yoneda : les objets représentables
forment une famille «~dense~» dans la catégorie des préfaisceaux, au sens
où tout préfaisceau s'obtient par colimites à partir des représentables.
\end{remarque}


%% ============================================================
\section{Exercices}\label{sec:exo-yoneda}
%% ============================================================

\begin{exercice}\label{exo:yoneda-groupes}
\index{Yoneda, lemme de!exercice}
Soit $G$ un groupe vu comme une catégorie $\mathbf{B}G$ à un seul objet
$*$ avec $\Hom(*, *) = G$.
\begin{enumerate}[label=(\alph*)]
\item Montrer qu'un préfaisceau sur $\mathbf{B}G$ est la même chose
    qu'un $G$-ensemble\index{G-ensemble@$G$-ensemble} (à droite).
\item Identifier l'image du plongement de Yoneda dans ce cas.
\item Énoncer et vérifier le lemme de Yoneda dans ce contexte :
    pour tout $G$-ensemble $X$,
    $\Nat(h_*, X) \cong X$ (en tant qu'ensembles).
\end{enumerate}
\end{exercice}

\begin{exercice}\label{exo:yoneda-preordre}
\index{préordre!et préfaisceaux}
Soit $(P, \leq)$ un ensemble préordonné vu comme une catégorie.
\begin{enumerate}[label=(\alph*)]
\item Montrer qu'un préfaisceau sur $P$ est un foncteur
    $P^{\op} \to \Set$, c'est-à-dire un ensemble $F(p)$ pour chaque
    $p \in P$ avec des applications de restriction
    $F(q) \to F(p)$ pour $p \leq q$, compatibles avec la composition.
\item Décrire le préfaisceau représentable $h_p$ pour $p \in P$.
\item Interpréter le lemme de Yoneda dans ce cas.
\end{enumerate}
\end{exercice}

\begin{exercice}\label{exo:representable-limites}
\index{foncteur représentable!et limites}
Montrer que tout foncteur représentable $h_A \colon \mathcal{C}^{\op} \to \Set$
préserve les limites. (C'est-à-dire, $h_A$ envoie les colimites de
$\mathcal{C}$ en limites de $\Set$.)
\end{exercice}

\begin{exercice}\label{exo:elements-comma}
\index{catégorie des éléments!et catégorie comma}
Montrer que la catégorie des éléments $\int_{\mathcal{C}} F$ est isomorphe
à la catégorie comma $(\mathbf{y} \downarrow F)$, où
$\mathbf{y} \colon \mathcal{C} \to \widehat{\mathcal{C}}$ est le plongement
de Yoneda.
\end{exercice}

\begin{exercice}\label{exo:coyoneda}
\index{Yoneda, lemme de!co-Yoneda}
\index{co-Yoneda, lemme de}
(\emph{Lemme de co-Yoneda.})
Soit $F \colon \mathcal{C} \to \Set$ un foncteur covariant.
Montrer que pour tout $A \in \mathcal{C}$,
\[
    F(A) \cong \int^{C \in \mathcal{C}} \Hom(C, A) \times F(C),
\]
où $\int^C$ désigne la coégalisatrice (coend)\index{coend}.
\end{exercice}

\begin{exercice}\label{exo:presheaf-complete}
\index{préfaisceau!complétude}
Soit $\mathcal{C}$ une catégorie petite. Montrer que
$\widehat{\mathcal{C}} = \Fun(\mathcal{C}^{\op}, \Set)$
admet toutes les petites limites et colimites, calculées
argument par argument.
\end{exercice}

\begin{exercice}\label{exo:yoneda-extension}
\index{extension de Kan!et Yoneda}
(\emph{Extension de Kan à gauche le long de Yoneda.})
Soit $\mathcal{C}$ une catégorie petite, $\mathcal{D}$ une catégorie
cocomplète, et $G \colon \mathcal{C} \to \mathcal{D}$ un foncteur.
Montrer que l'extension de Kan à gauche $\mathrm{Lan}_{\mathbf{y}} G
\colon \widehat{\mathcal{C}} \to \mathcal{D}$ existe et est donnée par
\[
    (\mathrm{Lan}_{\mathbf{y}} G)(F)
    = \colim_{(C, x) \in \int_{\mathcal{C}} F} G(C).
\]
\end{exercice}

\index{Yoneda, lemme de|)}
\index{préfaisceau|)}
\index{foncteur représentable|)}


%%%%%%%%%%%%%%%%%%%%%%%%%%%%%%%%%%%%%%%%%%%%%%%%%%%%%%%%%%%%
%%         CHAPITRE 6 : MONADES ET ALGÈBRES               %%
%%%%%%%%%%%%%%%%%%%%%%%%%%%%%%%%%%%%%%%%%%%%%%%%%%%%%%%%%%%%

\chapter{Monades et algèbres}\label{ch:monades}
\minitoc

\vspace{1em}

Les monades\index{monade|(} constituent un concept unificateur remarquable
qui relie la théorie des adjonctions, l'algèbre universelle et
l'informatique théorique.
Toute adjonction engendre une monade ; réciproquement, toute monade
provient d'une adjonction, et ce de deux manières canoniques
(Eilenberg--Moore et Kleisli).
Le théorème de Beck, résultat culminant de ce chapitre, caractérise
les catégories qui peuvent être décrites comme des «~catégories d'algèbres~»
sur une monade.

Ce chapitre suppose connue la théorie des adjonctions développée
aux chapitres précédents.

%% ============================================================
\section{Définition et premiers exemples}\label{sec:monades-def}
%% ============================================================

\begin{definition}[Monade]\label{def:monade}
\index{monade!définition}
\index{endofoncteur}
\index{multiplication!d'une monade}
\index{unité!d'une monade}
Une \textbf{monade} sur une catégorie $\mathcal{C}$ est un triplet
$(T, \eta, \mu)$ où :
\begin{itemize}
    \item $T \colon \mathcal{C} \to \mathcal{C}$ est un endofoncteur ;
    \item $\eta \colon \Id_{\mathcal{C}} \Rightarrow T$ est une
    transformation naturelle appelée \textbf{unité} ;
    \item $\mu \colon T^2 \Rightarrow T$ est une transformation naturelle
    appelée \textbf{multiplication},
\end{itemize}
satisfaisant les axiomes suivants :
\begin{enumerate}[label=(\roman*)]
\item \textbf{Associativité}\index{monade!associativité} :
    $\mu \circ T\mu = \mu \circ \mu T$, c'est-à-dire le diagramme
    suivant commute :
\[
\begin{tikzcd}
T^3 \arrow[r, "T\mu"] \arrow[d, "\mu T"'] & T^2 \arrow[d, "\mu"] \\
T^2 \arrow[r, "\mu"'] & T
\end{tikzcd}
\]

\item \textbf{Unités}\index{monade!axiomes d'unité} :
    $\mu \circ T\eta = \id_T = \mu \circ \eta T$, c'est-à-dire :
\[
\begin{tikzcd}
T \arrow[r, "T\eta"] \arrow[dr, equal] & T^2 \arrow[d, "\mu"] &
T \arrow[l, "\eta T"'] \arrow[dl, equal] \\
& T &
\end{tikzcd}
\]
\end{enumerate}
\end{definition}

\begin{remarque}\label{rem:monade-monoide}
\index{monade!comme monoïde}
\index{monoïde!dans une catégorie monoïdale}
Une monade sur $\mathcal{C}$ est exactement un monoïde dans la catégorie
monoïdale $(\mathrm{End}(\mathcal{C}), \circ, \Id_{\mathcal{C}})$
des endofoncteurs de $\mathcal{C}$, munie de la composition comme produit
tensoriel.
\end{remarque}

\begin{notation}\label{not:monade}
\index{monade!notation}
Par abus de langage, on note souvent la monade simplement $T$,
les transformations $\eta$ et $\mu$ étant sous-entendues.
On écrit parfois $\eta^T$ et $\mu^T$ lorsqu'il y a ambiguïté.
\end{notation}

\begin{definition}[Morphisme de monades]\label{def:morphisme-monades}
\index{monade!morphisme de monades}
Soient $(S, \eta^S, \mu^S)$ et $(T, \eta^T, \mu^T)$ deux monades
sur $\mathcal{C}$.
Un \textbf{morphisme de monades} $\sigma \colon S \to T$ est une
transformation naturelle $\sigma \colon S \Rightarrow T$ telle que :
\[
    \sigma \circ \eta^S = \eta^T, \qquad
    \sigma \circ \mu^S = \mu^T \circ (\sigma \star \sigma),
\]
où $\sigma \star \sigma = \sigma T \circ S\sigma = T\sigma \circ \sigma S
\colon S^2 \Rightarrow T^2$.
\end{definition}


%% ============================================================
\section{Monades issues d'adjonctions}\label{sec:monade-adjonction}
%% ============================================================

\begin{theoreme}\label{thm:adjonction-monade}
\index{adjonction!et monade}
\index{monade!issue d'une adjonction}
Toute adjonction $(F \dashv G, \eta, \varepsilon) \colon
\mathcal{C} \rightleftarrows \mathcal{D}$ engendre une monade
$(T, \eta, \mu)$ sur $\mathcal{C}$, où
\[
    T = G \circ F, \qquad
    \mu = G \varepsilon F \colon T^2 = GFGF \Rightarrow GF = T.
\]
\end{theoreme}

\begin{proof}
L'unité $\eta \colon \Id_{\mathcal{C}} \Rightarrow GF = T$ est
l'unité de l'adjonction. Il faut vérifier les axiomes de monade.

\emph{Associativité.}
On a $\mu \circ T\mu = G\varepsilon F \circ GF \cdot G\varepsilon F
= G\varepsilon F \circ G(F G\varepsilon F)
= G(\varepsilon F \circ FG\varepsilon F)
= G(\varepsilon \cdot \varepsilon FGF)$... Plus précisément,
$\mu \circ T\mu$ a pour composante en $C$ :
\[
    G(\varepsilon_{FC}) \circ GF(G(\varepsilon_{FC}))
    = G\bigl(\varepsilon_{FC} \circ F G(\varepsilon_{FC})\bigr).
\]
Et $\mu \circ \mu T$ a pour composante en $C$ :
\[
    G(\varepsilon_{FC}) \circ G(\varepsilon_{GFFC})
    = G\bigl(\varepsilon_{FC} \circ \varepsilon_{GFFC}\bigr).
\]
Or la naturalité de $\varepsilon$ appliquée au morphisme
$\varepsilon_{FC} \colon FGFC \to FC$ donne
$\varepsilon_{FC} \circ FG(\varepsilon_{FC})
= \varepsilon_{FC} \circ \varepsilon_{FGFC}$,
ce qui est exactement l'identité de gauche.

\emph{Unité à gauche.}
$\mu \circ \eta T$ a pour composante en $C$ :
$G(\varepsilon_{FC}) \circ \eta_{GFC}$.
Par l'identité triangulaire $G\varepsilon \circ \eta G = \id_G$,
en l'évaluant sur $FC$ on obtient
$G(\varepsilon_{FC}) \circ \eta_{GFC} = \id_{GFC} = \id_{TC}$.

\emph{Unité à droite.}
$\mu \circ T\eta$ a pour composante en $C$ :
$G(\varepsilon_{FC}) \circ GF(\eta_C)$.
Par l'identité triangulaire $\varepsilon F \circ F\eta = \id_F$,
en l'évaluant sur $C$ on obtient
$\varepsilon_{FC} \circ F(\eta_C) = \id_{FC}$,
d'où en appliquant $G$ :
$G(\varepsilon_{FC}) \circ GF(\eta_C) = G(\id_{FC}) = \id_{GFC}$.
\end{proof}

\begin{exemple}\label{ex:monades-adjonctions}
\index{monade!exemples}
\begin{enumerate}[label=(\roman*)]
\item \textbf{Monade libre sur les groupes.}%
\index{monade!libre!groupes}
    L'adjonction libre/oubli $F \dashv U \colon \Grp \rightleftarrows \Set$
    engendre la monade $T = UF$ sur $\Set$ : pour un ensemble $X$,
    $T(X)$ est l'ensemble sous-jacent au groupe libre $F(X)$.
    L'unité $\eta_X \colon X \to T(X)$ envoie chaque élément sur le
    mot de longueur $1$ correspondant. La multiplication
    $\mu_X \colon T^2(X) \to T(X)$ «~aplatit~» les mots de mots
    en les concaténant.

\item \textbf{Monade libre sur les modules.}%
\index{monade!libre!modules}
    Pour un anneau $R$, l'adjonction libre/oubli
    $F \dashv U \colon R\text{-}\mathsf{Mod} \rightleftarrows \Set$
    donne $T(X) = \bigoplus_{x \in X} R \cong R^{(X)}$.
    L'unité envoie $x$ sur le vecteur de base $e_x$
    et la multiplication est la $R$-linéarisation.

\item \textbf{Monade ensemble de puissance.}%
\index{monade!ensemble de puissance}
    Le foncteur $\mathcal{P} \colon \Set \to \Set$
    avec $\eta_X(x) = \{x\}$ et $\mu_X(\mathfrak{A}) = \bigcup \mathfrak{A}$
    pour $\mathfrak{A} \in \mathcal{P}(\mathcal{P}(X))$
    est une monade. Elle provient de l'auto-adjonction
    $\mathcal{P} \dashv \mathcal{P}^{\op}$ (à travers l'isomorphisme
    $\Set^{\op} \cong \mathsf{CABA}$).

\item \textbf{Monades en informatique.}%
\index{monade!en informatique}
    En programmation fonctionnelle (Haskell, Scala), les monades modélisent
    les effets de calcul. La monade \texttt{Maybe} correspond au foncteur
    $T(X) = X \sqcup \{*\}$ (valeur optionnelle). La monade \texttt{List}
    correspond au foncteur \emph{monoïde libre} $T(X) = \coprod_{n \geq 0} X^n$.
    La monade d'état $T(X) = (X \times S)^S$ pour un ensemble d'états $S$
    modélise les calculs avec état mutable.
\end{enumerate}
\end{exemple}

\begin{proposition}\label{prop:comonade}
\index{comonade}
\index{comonade!co-unité}
\index{comonade!comultiplication}
Dualement, toute adjonction $(F \dashv G)$ engendre une \textbf{comonade}
$(K, \varepsilon, \delta)$ sur $\mathcal{D}$, où
$K = F \circ G$, $\varepsilon$ est la co-unité de l'adjonction,
et $\delta = F\eta G \colon K \Rightarrow K^2$.
\end{proposition}

\begin{proof}
C'est le dual du théorème~\ref{thm:adjonction-monade}, en considérant
l'adjonction dans les catégories opposées.
\end{proof}


%% ============================================================
\section{Algèbres d'Eilenberg--Moore}\label{sec:eilenberg-moore}
%% ============================================================

La question réciproque se pose : toute monade provient-elle d'une adjonction ?
La réponse est oui, et il y a (au moins) deux constructions canoniques.
Nous commençons par celle d'Eilenberg--Moore\index{Eilenberg--Moore|(}.

\begin{definition}[$T$-algèbre]\label{def:T-algebre}
\index{T-algèbre@$T$-algèbre!définition}
\index{algèbre!d'une monade}
\index{application de structure}
Soit $(T, \eta, \mu)$ une monade sur $\mathcal{C}$.
Une \textbf{$T$-algèbre} (ou \textbf{algèbre d'Eilenberg--Moore})
est un couple $(A, a)$ où :
\begin{itemize}
    \item $A \in \Ob(\mathcal{C})$ est un objet (le \textbf{porteur}) ;
    \item $a \colon T(A) \to A$ est un morphisme
    (l'\textbf{application de structure}),
\end{itemize}
satisfaisant :
\begin{enumerate}[label=(\roman*)]
\item \textbf{Unité}\index{T-algèbre@$T$-algèbre!axiome d'unité} :
    $a \circ \eta_A = \id_A$ :
\[
\begin{tikzcd}
A \arrow[r, "\eta_A"] \arrow[dr, equal] & T(A) \arrow[d, "a"] \\
& A
\end{tikzcd}
\]

\item \textbf{Associativité}\index{T-algèbre@$T$-algèbre!axiome d'associativité} :
    $a \circ \mu_A = a \circ T(a)$ :
\[
\begin{tikzcd}
T^2(A) \arrow[r, "T(a)"] \arrow[d, "\mu_A"'] & T(A) \arrow[d, "a"] \\
T(A) \arrow[r, "a"'] & A
\end{tikzcd}
\]
\end{enumerate}
\end{definition}

\begin{definition}[Morphisme de $T$-algèbres]\label{def:morphisme-T-algebre}
\index{T-algèbre@$T$-algèbre!morphisme}
Un \textbf{morphisme de $T$-algèbres} de $(A, a)$ vers $(B, b)$ est
un morphisme $f \colon A \to B$ dans $\mathcal{C}$ tel que
$f \circ a = b \circ T(f)$ :
\[
\begin{tikzcd}
T(A) \arrow[r, "T(f)"] \arrow[d, "a"'] & T(B) \arrow[d, "b"] \\
A \arrow[r, "f"'] & B
\end{tikzcd}
\]
\end{definition}

\begin{definition}[Catégorie d'Eilenberg--Moore]\label{def:categorie-EM}
\index{catégorie!d'Eilenberg--Moore}
\index{Eilenberg--Moore!catégorie}
\index{$\mathcal{C}^T$}
La \textbf{catégorie d'Eilenberg--Moore} de la monade $T$,
notée $\mathcal{C}^T$, est la catégorie dont les objets sont les
$T$-algèbres et les morphismes sont les morphismes de $T$-algèbres.
La composition et les identités sont héritées de $\mathcal{C}$.
\end{definition}

\begin{exemple}\label{ex:T-algebres}
\index{T-algèbre@$T$-algèbre!exemples}
\begin{enumerate}[label=(\roman*)]
\item \textbf{Monade libre sur les groupes.}
    La monade $T = UF$ sur $\Set$ (foncteur mot).
    Une $T$-algèbre est un ensemble $X$ muni d'une application
    $a \colon T(X) \to X$ qui est associative et unitale.
    Cela revient exactement à munir $X$ d'une structure de groupe.
    Ainsi $\Set^T \simeq \Grp$.\index{monade!libre!algèbres}

\item \textbf{Monade ensemble de puissance.}
    Les $\mathcal{P}$-algèbres sur $\Set$ sont les treillis complets
    avec morphismes préservant les suprema arbitraires.
    \index{treillis complet}

\item \textbf{Algèbre libre.}
    Pour toute monade $T$ et tout objet $C \in \mathcal{C}$,
    le couple $(T(C), \mu_C)$ est une $T$-algèbre, appelée
    la \textbf{$T$-algèbre libre}\index{T-algèbre@$T$-algèbre!libre}
    engendrée par $C$.
    En effet, $\mu_C \circ \eta_{TC} = \id_{TC}$ (unité de droite)
    et $\mu_C \circ \mu_{TC} = \mu_C \circ T\mu_C$ (associativité).
\end{enumerate}
\end{exemple}

\begin{theoreme}\label{thm:adjonction-EM}
\index{adjonction!d'Eilenberg--Moore}
\index{Eilenberg--Moore!adjonction}
Soit $(T, \eta, \mu)$ une monade sur $\mathcal{C}$.
Il existe une adjonction
\[
    F^T \dashv U^T \colon \mathcal{C}^T \rightleftarrows \mathcal{C}
\]
telle que la monade associée est $(T, \eta, \mu)$, où :
\begin{itemize}
    \item le foncteur d'oubli $U^T \colon \mathcal{C}^T \to \mathcal{C}$
    est défini par $U^T(A, a) = A$ et $U^T(f) = f$ ;
    \item le foncteur libre $F^T \colon \mathcal{C} \to \mathcal{C}^T$
    est défini par $F^T(C) = (TC, \mu_C)$ et
    $F^T(f) = T(f)$.
\end{itemize}
\end{theoreme}

\begin{proof}
\emph{$F^T$ est bien défini.}
Pour tout $C \in \mathcal{C}$, $(TC, \mu_C)$ est une $T$-algèbre
libre (exemple~\ref{ex:T-algebres}(iii)).
Pour un morphisme $f \colon C \to D$, $T(f) \colon TC \to TD$ est
un morphisme de $T$-algèbres car
$T(f) \circ \mu_C = \mu_D \circ T^2(f)$ par naturalité de $\mu$.

\emph{Adjonction.}
On définit l'unité $\eta^{\mathrm{adj}}_C = \eta_C \colon C \to U^T F^T(C) = TC$
et la co-unité $\varepsilon^{\mathrm{adj}}_{(A,a)} = a \colon F^T U^T(A,a)
= (TA, \mu_A) \to (A,a)$.
Vérifions que $a$ est un morphisme de $T$-algèbres :
$a \circ \mu_A = a \circ Ta$ par l'axiome d'associativité de $(A,a)$.

\emph{Identités triangulaires.}
\begin{itemize}
    \item $U^T \varepsilon \circ \eta_{U^T} \colon U^T \Rightarrow U^T$:
    pour $(A,a)$, la composante est $a \circ \eta_A = \id_A$
    par l'axiome d'unité. Donc c'est $\id_{U^T}$.

    \item $\varepsilon_{F^T} \circ F^T\eta \colon F^T \Rightarrow F^T$:
    pour $C$, la composante est $\mu_C \circ T\eta_C = \id_{TC}$
    par l'axiome d'unité de droite de la monade. Donc c'est $\id_{F^T}$.
\end{itemize}

\emph{La monade associée est $T$.}
$U^T F^T(C) = TC$, et la multiplication est
$U^T \varepsilon_{F^T C} = \mu_C$. L'unité est $\eta_C$.
\end{proof}


%% ============================================================
\section{Catégorie de Kleisli}\label{sec:kleisli}
%% ============================================================

\begin{definition}[Catégorie de Kleisli]\label{def:categorie-kleisli}
\index{Kleisli!catégorie de}
\index{catégorie!de Kleisli}
\index{$\mathcal{C}_T$}
Soit $(T, \eta, \mu)$ une monade sur $\mathcal{C}$.
La \textbf{catégorie de Kleisli} $\mathcal{C}_T$ est définie par :
\begin{itemize}
    \item $\Ob(\mathcal{C}_T) = \Ob(\mathcal{C})$ ;
    \item pour $A, B \in \Ob(\mathcal{C}_T)$, on pose
    $\Hom_{\mathcal{C}_T}(A, B) = \Hom_{\mathcal{C}}(A, TB)$ ;
    \item la composition de $f \colon A \to TB$ et $g \colon B \to TC$
    dans $\mathcal{C}_T$ est
    \[
        g \circ_T f \coloneqq \mu_C \circ Tg \circ f \colon A \to TC ;
    \]
    \item l'identité de $A$ dans $\mathcal{C}_T$ est $\eta_A \colon A \to TA$.
\end{itemize}
\end{definition}

\begin{proposition}\label{prop:kleisli-categorie}
\index{Kleisli!bonne définition}
La catégorie de Kleisli $\mathcal{C}_T$ est bien une catégorie.
\end{proposition}

\begin{proof}
\emph{Associativité.}
Pour $f \colon A \to TB$, $g \colon B \to TC$, $h \colon C \to TD$,
\begin{align*}
    h \circ_T (g \circ_T f)
    &= \mu_D \circ Th \circ (\mu_C \circ Tg \circ f) \\
    &= \mu_D \circ Th \circ \mu_C \circ Tg \circ f \\
    &= \mu_D \circ \mu_{TD} \circ T^2 h \circ Tg \circ f
    \quad \text{(naturalité de $\mu$)}.
\end{align*}
D'autre part,
\begin{align*}
    (h \circ_T g) \circ_T f
    &= \mu_D \circ T(\mu_D \circ Th \circ g) \circ f \\
    &= \mu_D \circ T\mu_D \circ T^2 h \circ Tg \circ f.
\end{align*}
Par l'associativité de la monade $\mu_D \circ \mu_{TD} = \mu_D \circ T\mu_D$,
les deux expressions sont égales.

\emph{Unité à gauche.}
$f \circ_T \eta_A = \mu_B \circ Tf \circ \eta_A
= \mu_B \circ \eta_{TB} \circ f = f$
par naturalité de $\eta$ et l'axiome d'unité.

\emph{Unité à droite.}
$\eta_B \circ_T f = \mu_B \circ T\eta_B \circ f = \id_{TB} \circ f = f$
par l'axiome d'unité de la monade.
\end{proof}

\begin{theoreme}\label{thm:adjonction-kleisli}
\index{Kleisli!adjonction}
\index{adjonction!de Kleisli}
Il existe une adjonction $F_T \dashv G_T \colon \mathcal{C}_T
\rightleftarrows \mathcal{C}$ telle que la monade associée est $T$, où :
\begin{itemize}
    \item $F_T \colon \mathcal{C} \to \mathcal{C}_T$ est défini par
    $F_T(A) = A$ sur les objets et $F_T(f) = \eta_B \circ f$ pour
    $f \colon A \to B$ ;
    \item $G_T \colon \mathcal{C}_T \to \mathcal{C}$ est défini par
    $G_T(A) = TA$ et $G_T(f) = \mu_B \circ Tf$ pour $f \colon A \to TB$.
\end{itemize}
\end{theoreme}

\begin{proof}
La vérification est directe. L'unité de l'adjonction est
$\eta_A \colon A \to G_T F_T(A) = TA$, et la co-unité est
$\varepsilon_A = \id_{TA} \colon F_T G_T(A) = TA \to TA$ dans
$\mathcal{C}_T$ (vue comme $\id_{TA} \in \Hom(TA, TA)$).
Les identités triangulaires se vérifient à l'aide des axiomes de monade.
\end{proof}

\begin{remarque}\label{rem:kleisli-EM-universalite}
\index{Kleisli!propriété universelle}
\index{Eilenberg--Moore!propriété universelle}
La catégorie de Kleisli est \emph{initiale} et celle d'Eilenberg--Moore
est \emph{terminale} parmi les adjonctions engendrant la monade $T$.
Plus précisément, si $(F \dashv G) \colon \mathcal{C} \rightleftarrows
\mathcal{D}$ est une adjonction engendrant $T$, il existe des foncteurs
uniques (à isomorphisme unique près) formant le diagramme commutatif :
\[
\begin{tikzcd}
& \mathcal{D} \arrow[dd, "K"] \\
\mathcal{C}_T \arrow[ur, "J"] \arrow[dr] & \\
& \mathcal{C}^T
\end{tikzcd}
\]
où $J$ est pleinement fidèle et $K$ est le foncteur de comparaison.
\end{remarque}


%% ============================================================
\section{Théorème de monadicité de Beck}\label{sec:beck}
%% ============================================================

Le théorème de Beck\index{Beck, théorème de|(} caractérise les adjonctions
pour lesquelles le foncteur de comparaison est une équivalence.

\begin{definition}[Paire coégalisatrice scindée]\label{def:paire-coeg-scindee}
\index{paire coégalisatrice scindée}
\index{coégalisateur!scindé}
Soit $\mathcal{C}$ une catégorie.
Une \textbf{paire coégalisatrice scindée} est un diagramme
\[
\begin{tikzcd}
A \arrow[r, shift left=0.5ex, "f"] \arrow[r, shift right=0.5ex, "g"'] &
B \arrow[r, "e"] \arrow[l, bend right=40, "s"'] &
C \arrow[l, bend right=40, "t"']
\end{tikzcd}
\]
tel que $e \circ f = e \circ g$, $e \circ t = \id_C$,
$f \circ s = \id_B$ et $g \circ s = t \circ e$.
\end{definition}

\begin{lemme}\label{lem:coeg-scindee-absolue}
\index{coégalisateur!absolu}
\index{paire coégalisatrice scindée!est coégalisateur}
Toute paire coégalisatrice scindée est un coégalisateur absolu,
c'est-à-dire que $e$ est le coégalisateur de $(f, g)$ et ce coégalisateur
est préservé par tout foncteur.
\end{lemme}

\begin{proof}
Montrons que $e$ est le coégalisateur de $(f, g)$.
On sait que $e \circ f = e \circ g$.
Soit $h \colon B \to D$ tel que $h \circ f = h \circ g$.
Posons $\bar{h} = h \circ t \colon C \to D$. Alors
$\bar{h} \circ e = h \circ t \circ e = h \circ g \circ s \cdot$...
Plus précisément :
\[
    \bar{h} \circ e = h \circ t \circ e.
\]
Or $g \circ s = t \circ e$, donc $t \circ e = g \circ s$ et
\[
    h \circ t \circ e = h \circ g \circ s = h \circ f \circ s = h \circ \id_B = h.
\]
Donc $\bar{h} \circ e = h$.

Pour l'unicité, si $\bar{h}' \circ e = h$, alors
$\bar{h}' = \bar{h}' \circ e \circ t = h \circ t = \bar{h}$.

Puisque les identités scindantes se transportent par tout foncteur $U$
(en appliquant $U$ aux morphismes $s, t$),
le coégalisateur est préservé par $U$.
\end{proof}

\begin{definition}[Foncteur de comparaison]\label{def:foncteur-comparaison}
\index{foncteur!de comparaison}
\index{comparaison, foncteur de}
Soit $(F \dashv G, \eta, \varepsilon)$ une adjonction
$\mathcal{C} \rightleftarrows \mathcal{D}$ engendrant la monade $T = GF$.
Le \textbf{foncteur de comparaison} est le foncteur
$K \colon \mathcal{D} \to \mathcal{C}^T$ défini par :
\begin{itemize}
    \item $K(D) = (GD,\, G\varepsilon_D)$ pour tout $D \in \mathcal{D}$ ;
    \item $K(f) = Gf$ pour tout morphisme $f$ dans $\mathcal{D}$.
\end{itemize}
\end{definition}

\begin{proposition}\label{prop:comparaison-bien-defini}
\index{foncteur!de comparaison!bonne définition}
Le foncteur de comparaison $K$ est bien défini :
$(GD, G\varepsilon_D)$ est une $T$-algèbre pour tout $D$.
\end{proposition}

\begin{proof}
\emph{Unité.}
$G\varepsilon_D \circ \eta_{GD} = \id_{GD}$ par l'identité triangulaire.

\emph{Associativité.}
$G\varepsilon_D \circ \mu_{GD} = G\varepsilon_D \circ G\varepsilon_{FGD}
= G(\varepsilon_D \circ \varepsilon_{FGD})$.
Par ailleurs, $G\varepsilon_D \circ GF(G\varepsilon_D)
= G(\varepsilon_D \circ FG\varepsilon_D)$.
La naturalité de $\varepsilon$ appliquée à $\varepsilon_D$ donne
$\varepsilon_D \circ FG\varepsilon_D = \varepsilon_D \circ \varepsilon_{FGD}$.
\end{proof}

\begin{definition}[Foncteur monadique]\label{def:foncteur-monadique}
\index{foncteur!monadique}
\index{monadique, foncteur}
Un foncteur $G \colon \mathcal{D} \to \mathcal{C}$ est dit
\textbf{monadique} s'il admet un adjoint à gauche $F$ et si le
foncteur de comparaison $K \colon \mathcal{D} \to \mathcal{C}^T$
(pour la monade $T = GF$) est une équivalence de catégories.
Il est \textbf{strictement monadique} si $K$ est un isomorphisme
de catégories.
\end{definition}

\begin{definition}[$G$-paire coégalisatrice scindée]\label{def:G-paire-coeg}
\index{paire coégalisatrice!$G$-scindée}
Soit $G \colon \mathcal{D} \to \mathcal{C}$ un foncteur.
Une paire de morphismes $f, g \colon A \rightrightarrows B$ dans $\mathcal{D}$
est une \textbf{$G$-paire coégalisatrice scindée} si la paire
$Gf, Gg \colon GA \rightrightarrows GB$ admet une paire coégalisatrice
scindée dans $\mathcal{C}$.
\end{definition}

\begin{theoreme}[Théorème de monadicité de Beck]\label{thm:beck}
\index{Beck, théorème de!énoncé}
\index{monadicité, théorème de}
Soit $G \colon \mathcal{D} \to \mathcal{C}$ un foncteur admettant
un adjoint à gauche $F$. Alors $G$ est monadique si et seulement si :
\begin{enumerate}[label=(\roman*)]
\item $G$ reflète les isomorphismes\index{foncteur!qui reflète les isomorphismes}
    (c'est-à-dire : si $Gf$ est un isomorphisme dans $\mathcal{C}$,
    alors $f$ est un isomorphisme dans $\mathcal{D}$) ;
\item $\mathcal{D}$ admet les coégalisateurs de toutes les
    $G$-paires coégalisatrices scindées, et $G$ les préserve.
\end{enumerate}
\end{theoreme}

\begin{proof}
\index{Beck, théorème de!preuve}
La preuve se décompose en deux implications.

\medskip
\textbf{($\Rightarrow$) Condition nécessaire.}

Supposons $G$ monadique, c'est-à-dire que $K \colon \mathcal{D} \to
\mathcal{C}^T$ est une équivalence.

\emph{(i) $G$ reflète les isomorphismes.}
Puisque $U^T \circ K = G$ et que $K$ est une équivalence,
il suffit de montrer que $U^T$ reflète les isomorphismes.
Soit $f \colon (A, a) \to (B, b)$ un morphisme de $T$-algèbres
tel que $U^T(f) = f \colon A \to B$ est un isomorphisme dans $\mathcal{C}$.
Alors $f^{-1}$ existe dans $\mathcal{C}$, et il faut montrer que
$f^{-1}$ est un morphisme de $T$-algèbres, c'est-à-dire
$f^{-1} \circ b = a \circ T(f^{-1})$.
Or $f \circ a = b \circ Tf$ donne
$a = f^{-1} \circ b \circ Tf$, donc
$a \circ T(f^{-1}) = f^{-1} \circ b \circ Tf \circ T(f^{-1})
= f^{-1} \circ b$.

\emph{(ii) Existence et préservation des coégalisateurs.}
Puisque $K$ est une équivalence, il suffit de montrer le résultat
pour $\mathcal{C}^T$. Si $(f, g) \colon (A,a) \rightrightarrows (B,b)$
est une paire dans $\mathcal{C}^T$ telle que $U^T f, U^T g$ admet
une paire coégalisatrice scindée dans $\mathcal{C}$, alors le coégalisateur
scindé est absolu (lemme~\ref{lem:coeg-scindee-absolue}) et on peut le relever
en une $T$-algèbre.

\medskip
\textbf{($\Leftarrow$) Condition suffisante.}

Supposons (i) et (ii). Montrons que $K$ est une équivalence
en vérifiant qu'il est pleinement fidèle et essentiellement surjectif.

\emph{$K$ est fidèle.}
$K$ est fidèle car $U^T \circ K = G$ et $U^T$ est fidèle
(comme foncteur d'oubli d'algèbres).

\emph{$K$ est plein.}
Soit $\varphi \colon K(D) \to K(D')$ un morphisme de $T$-algèbres,
c'est-à-dire $\varphi \colon GD \to GD'$ avec
$\varphi \circ G\varepsilon_D = G\varepsilon_{D'} \circ GF\varphi$.
Considérons le diagramme dans $\mathcal{D}$ :
\[
\begin{tikzcd}
FGFGD \arrow[r, shift left, "FG\varepsilon_D"] \arrow[r, shift right, "\varepsilon_{FGD}"'] &
FGD \arrow[r, "\varepsilon_D"] & D
\end{tikzcd}
\]
En appliquant $G$, on obtient une paire coégalisatrice scindée
(les sections proviennent de $\eta$). Par (ii), $\varepsilon_D$ est
le coégalisateur de cette paire dans $\mathcal{D}$.
De même pour $D'$. La condition de morphisme de $T$-algèbres sur
$\varphi$ assure que $F\varphi \colon FGD \to FGD'$ est compatible,
et par la propriété universelle du coégalisateur, il existe un unique
$f \colon D \to D'$ avec $Gf = \varphi$. Donc $K(f) = \varphi$.

\emph{$K$ est essentiellement surjectif.}
Soit $(A, a)$ une $T$-algèbre. Considérons la paire
\[
\begin{tikzcd}
FTA \arrow[r, shift left, "Fa"] \arrow[r, shift right, "\varepsilon_{FA} \circ FT(\eta_A)"'] &
FA
\end{tikzcd}
\]
Plus précisément, on considère les deux morphismes
$F(a), \varepsilon_{FA} \colon FGFA = FTA \rightrightarrows FA$.
En appliquant $G$, on obtient $GFa, G\varepsilon_{FA} \colon
GFGFA \rightrightarrows GFA$, qui est une $G$-paire coégalisatrice
scindée (les sections sont données par $\eta_{GFA}$ et $GF\eta_A$,
en utilisant les axiomes de la $T$-algèbre).

Par (ii), soit $e \colon FA \to D$ le coégalisateur de
$(F(a), \varepsilon_{FA})$ dans $\mathcal{D}$, préservé par $G$.
Alors $Ge \colon GFA = TA \to GD$ est le coégalisateur de
$(GFa, G\varepsilon_{FA}) = (Ta, \mu_A)$ dans $\mathcal{C}$.
Or $a$ coégalise cette paire ($a \circ Ta = a \circ \mu_A$
par l'axiome d'associativité), donc il existe un unique
$\bar{a} \colon GD \to A$ avec $\bar{a} \circ Ge = a$.

On vérifie que $\bar{a}$ est un isomorphisme.
En effet, $\bar{a} \circ Ge \circ \eta_A = a \circ \eta_A = \id_A$,
et on montre que $Ge \circ \eta_A \circ \bar{a} = \id_{GD}$
en utilisant que $Ge$ est un épimorphisme (étant un coégalisateur).
Donc $G(\bar{a})$... plus précisément, on montre que $\bar{a}$ induit
un isomorphisme $GD \cong A$ et par (i) ($G$ reflète les isomorphismes),
l'isomorphisme se relève.

Ainsi $(A, a) \cong K(D) = (GD, G\varepsilon_D)$ dans $\mathcal{C}^T$.
\end{proof}

\begin{remarque}\label{rem:beck-versions}
\index{Beck, théorème de!version stricte}
Il existe une version «~stricte~» du théorème de Beck :
$G$ est \emph{strictement} monadique si et seulement si
$G$ reflète les isomorphismes et $\mathcal{D}$ admet les
coégalisateurs de \emph{toutes} les paires réflexives
(c'est-à-dire $f, g \colon A \rightrightarrows B$ avec une section
commune $s \colon B \to A$, $fs = gs = \id_B$), et $G$ les préserve.

On utilise aussi fréquemment le \emph{théorème de monadicité par
coégalisateurs réflexifs}\index{coégalisateur!réflexif} :
si $G$ crée les coégalisateurs de $G$-paires coégalisatrices scindées
(au lieu de simplement les préserver), alors $G$ est strictement monadique.
\end{remarque}

\index{Beck, théorème de|)}


%% ============================================================
\section{Applications : catégories algébriques}\label{sec:categories-algebriques}
%% ============================================================

Le théorème de Beck permet de reconnaître les catégories
«~algébriques~», c'est-à-dire celles qui sont monadiques sur $\Set$.

\begin{definition}[Catégorie algébrique]\label{def:categorie-algebrique}
\index{catégorie!algébrique}
Une catégorie $\mathcal{D}$ est dite \textbf{algébrique} (au sens de
la monadicité) si le foncteur d'oubli $G \colon \mathcal{D} \to \Set$
est monadique.
\end{definition}

\begin{corollaire}\label{cor:grp-monadique}
\index{Grp@$\Grp$!monadicité}
\index{Ab@$\Ab$!monadicité}
\index{Ring@$\mathsf{Ring}$!monadicité}
Les catégories $\Grp$, $\Ab$, $R\text{-}\mathsf{Mod}$,
$\mathsf{Ring}$ sont monadiques sur $\Set$.
\end{corollaire}

\begin{proof}
Nous vérifions les conditions du théorème de Beck pour le foncteur
d'oubli $U \colon \Grp \to \Set$ (les autres cas sont analogues).

\emph{$U$ admet un adjoint à gauche.}
C'est le foncteur groupe libre $F \colon \Set \to \Grp$.

\emph{$U$ reflète les isomorphismes.}
Un morphisme de groupes $\varphi \colon G \to H$ tel que $U\varphi$
est bijectif est un isomorphisme de groupes (l'inverse ensembliste
est automatiquement un morphisme de groupes :
$\varphi^{-1}(xy) = \varphi^{-1}(x)\varphi^{-1}(y)$
car $\varphi$ est un morphisme bijectif).

\emph{$\Grp$ admet les coégalisateurs des $U$-paires coégalisatrices
scindées et $U$ les préserve.}
Plus généralement, $\Grp$ admet tous les coégalisateurs
(c'est la catégorie quotient par la relation d'équivalence engendrée).
Si la paire est $U$-scindée, le coégalisateur dans $\Grp$
est préservé par $U$ car les paires coégalisatrices scindées
donnent des coégalisateurs absolus
(lemme~\ref{lem:coeg-scindee-absolue}).
\end{proof}

\begin{exemple}\label{ex:non-monadique}
\index{catégorie!non monadique}
\begin{enumerate}[label=(\roman*)]
\item \textbf{$\Top$ n'est pas monadique sur $\Set$.}
    Le foncteur d'oubli $U \colon \Top \to \Set$ ne reflète pas
    les isomorphismes : une bijection continue n'est pas nécessairement
    un homéomorphisme. Donc $\Top$ n'est pas monadique sur $\Set$.

\item \textbf{Les espaces de Hausdorff compacts.}\index{espace de Hausdorff compact}
    Le foncteur d'oubli $U \colon \mathsf{CHaus} \to \Set$
    est monadique (théorème de Manes\index{Manes, théorème de}).
    La monade associée est la monade ultrafiltres
    $\beta \colon \Set \to \Set$\index{monade!ultrafiltres},
    qui envoie un ensemble sur l'ensemble de ses ultrafiltres.
\end{enumerate}
\end{exemple}

\begin{proposition}\label{prop:monadique-proprietes}
\index{catégorie!monadique!propriétés}
Soit $G \colon \mathcal{D} \to \mathcal{C}$ un foncteur monadique.
Alors :
\begin{enumerate}[label=(\roman*)]
\item $G$ reflète les isomorphismes ;
\item $G$ crée les limites que $\mathcal{C}$ possède
    (en particulier $\mathcal{D}$ est complète si $\mathcal{C}$ l'est) ;
\item $G$ est fidèle et conservatif.
\end{enumerate}
\end{proposition}

\begin{proof}
(i) est le théorème de Beck. Pour (ii), on montre que $U^T$ crée
les limites. Soit $\{(A_i, a_i)\}$ un diagramme de $T$-algèbres et
$L = \lim_i A_i$ dans $\mathcal{C}$ avec projections $\pi_i$.
Alors $TL = T(\lim_i A_i)$ et on définit $\ell \colon TL \to L$
comme l'unique morphisme tel que $\pi_i \circ \ell = a_i \circ T\pi_i$
pour tout $i$. On vérifie que $(L, \ell)$ est une $T$-algèbre et
que c'est la limite dans $\mathcal{C}^T$. Comme $K$ est une équivalence,
le résultat se transporte à $\mathcal{D}$.

(iii) découle de (i) et de l'existence de l'adjoint à gauche.
\end{proof}

\begin{theoreme}[Critère de monadicité de Beck simplifié]%
\label{thm:beck-simplifie}
\index{Beck, théorème de!critère simplifié}
\index{monadicité!critère simplifié}
Soit $G \colon \mathcal{D} \to \mathcal{C}$ un foncteur admettant un
adjoint à gauche. Si $G$ reflète les isomorphismes et si $\mathcal{D}$
admet tous les coégalisateurs réflexifs\index{coégalisateur!réflexif}
et $G$ les préserve, alors $G$ est monadique.
\end{theoreme}

\begin{proof}
Toute $G$-paire coégalisatrice scindée est en particulier réflexive
(la section $s$ fournit la section commune). Donc les hypothèses du
théorème de Beck~\ref{thm:beck} sont satisfaites.
\end{proof}


%% ============================================================
\section{Exercices}\label{sec:exo-monades}
%% ============================================================

\begin{exercice}\label{exo:monade-idempotente}
\index{monade!idempotente}
Une monade $(T, \eta, \mu)$ est dite \textbf{idempotente} si
$\mu \colon T^2 \Rightarrow T$ est un isomorphisme naturel.
\begin{enumerate}[label=(\alph*)]
\item Montrer que $T$ est idempotente si et seulement si
    $T\eta = \eta T \colon T \Rightarrow T^2$.
\item Montrer que $T$ est idempotente si et seulement si
    $U^T \colon \mathcal{C}^T \to \mathcal{C}$ est pleinement fidèle.
\item En déduire que les $T$-algèbres correspondent à une
    sous-catégorie pleine réflexive de $\mathcal{C}$.
\end{enumerate}
\end{exercice}

\begin{exercice}\label{exo:kleisli-informatique}
\index{Kleisli!en informatique}
Soit $T$ la monade \texttt{Maybe} sur $\Set$ :
$T(X) = X \sqcup \{\bot\}$, $\eta_X(x) = x$,
$\mu_X(x) = x$ si $x \in X$ et $\mu_X(\bot) = \bot$.
\begin{enumerate}[label=(\alph*)]
\item Vérifier les axiomes de monade.
\item Décrire la catégorie de Kleisli $\Set_T$ : quels sont les morphismes
    de $A$ vers $B$ ? Interpréter comme des fonctions partielles.
\item Décrire les $T$-algèbres (catégorie d'Eilenberg--Moore).
\end{enumerate}
\end{exercice}

\begin{exercice}\label{exo:monade-distribution}
\index{monade!de distribution}
Soit $T \colon \Set \to \Set$ le foncteur qui envoie $X$ sur l'ensemble
des distributions de probabilité finiment supportées sur $X$ :
\[
    T(X) = \Bigl\{ p \colon X \to [0,1] \;\Big|\;
    \mathrm{supp}(p) \text{ fini}, \sum_{x \in X} p(x) = 1 \Bigr\}.
\]
\begin{enumerate}[label=(\alph*)]
\item Définir $\eta$ et $\mu$ pour en faire une monade.
\item Identifier les $T$-algèbres comme des espaces convexes
    (ensembles munis de combinaisons convexes).
\end{enumerate}
\end{exercice}

\begin{exercice}\label{exo:adjonction-monade-aller-retour}
\index{adjonction!et monade}
Soit $(F \dashv G, \eta, \varepsilon)$ une adjonction engendrant
la monade $T$ et la comonade $K$.
\begin{enumerate}[label=(\alph*)]
\item Montrer que les identités triangulaires de l'adjonction sont
    exactement ce qu'il faut pour que $\eta$ et $\mu = G\varepsilon F$
    satisfassent les axiomes de monade.
\item Montrer que si $F$ est pleinement fidèle, alors
    $\eta$ est un isomorphisme naturel, et $T$ est idempotente.
\item Montrer que si $G$ est pleinement fidèle, alors
    $\varepsilon$ est un isomorphisme naturel, et $K$ est idempotente.
\end{enumerate}
\end{exercice}

\begin{exercice}\label{exo:algebres-libres-colimites}
\index{T-algèbre@$T$-algèbre!libre!colimites}
Soit $T$ une monade sur $\mathcal{C}$.
\begin{enumerate}[label=(\alph*)]
\item Montrer que le foncteur libre $F^T \colon \mathcal{C} \to \mathcal{C}^T$
    préserve les colimites.
\item Montrer que toute $T$-algèbre $(A, a)$ est le coégalisateur dans
    $\mathcal{C}^T$ de la paire
    \[
    \begin{tikzcd}
    F^T T(A) \arrow[r, shift left, "F^T(a)"]
    \arrow[r, shift right, "\varepsilon_{F^T(A)}"'] & F^T(A)
    \end{tikzcd}
    \]
    c'est-à-dire que toute algèbre est un coégalisateur d'algèbres libres.
\end{enumerate}
\end{exercice}

\begin{exercice}\label{exo:beck-top}
\index{Top@$\Top$!non monadicité}
Donner un exemple explicite montrant que le foncteur d'oubli
$U \colon \Top \to \Set$ ne reflète pas les isomorphismes.
En déduire que $\Top$ n'est pas monadique sur $\Set$.
\end{exercice}

\begin{exercice}\label{exo:distributive-law}
\index{loi distributive}
Soient $S$ et $T$ deux monades sur $\mathcal{C}$.
Une \textbf{loi distributive} de $S$ sur $T$ est une
transformation naturelle $\lambda \colon TS \Rightarrow ST$ compatible
avec les unités et multiplications des deux monades.
\begin{enumerate}[label=(\alph*)]
\item Écrire les quatre diagrammes de compatibilité que $\lambda$ doit
    satisfaire.
\item Montrer qu'une loi distributive $\lambda$ définit une monade
    $ST$ sur $\mathcal{C}$ avec unité $\eta^S \eta^T$ et multiplication
    $\mu^S \mu^T \circ S\lambda T$.
\item Donner un exemple avec $S = (-) \times M$ (monade écrivain pour
    un monoïde $M$) et $T = (-)^R$ (monade lecteur pour un ensemble $R$).
\end{enumerate}
\end{exercice}

\begin{exercice}\label{exo:monade-codensity}
\index{monade!de codensité}
(\emph{Monade de codensité.})
Soit $G \colon \mathcal{D} \to \mathcal{C}$ un foncteur qui
n'admet pas nécessairement d'adjoint à gauche.
On définit l'extension de Kan à droite
$\mathrm{Ran}_G G \colon \mathcal{C} \to \mathcal{C}$ (lorsqu'elle existe).
\begin{enumerate}[label=(\alph*)]
\item Montrer que $T = \mathrm{Ran}_G G$ est naturellement muni d'une
    structure de monade, appelée la \textbf{monade de codensité} de $G$.
\item Montrer que lorsque $G$ admet un adjoint à gauche $F$,
    la monade de codensité coïncide avec $GF$.
\item Pour $G \colon \mathsf{FinSet} \hookrightarrow \Set$ l'inclusion,
    montrer que la monade de codensité est la monade des ultrafiltres $\beta$.
\end{enumerate}
\end{exercice}

\index{Eilenberg--Moore|)}
\index{monade|)}
